
\documentclass[11pt]{article}
\usepackage{amsmath,amssymb,amsthm,mathrsfs}
\usepackage[a4paper,margin=1in]{geometry}

\title{Viscous Hamilton--Jacobi Flow, Hessian Evolution, and Riccati Bounds}
\author{Synthetic Compilation of Established Results}
\date{\today}

\newtheorem{theorem}{Theorem}[section]
\newtheorem{lemma}[theorem]{Lemma}
\newtheorem{proposition}[theorem]{Proposition}
\newtheorem{corollary}[theorem]{Corollary}
\newtheorem{definition}[theorem]{Definition}
\newtheorem{remark}[theorem]{Remark}

\begin{document}
\maketitle

\begin{abstract}
We derive in a finite--dimensional setting the viscous
Hamilton--Jacobi (vHJ) equation for an effective action evolving under
heat flow, the associated matrix--valued Hessian reaction--diffusion
equation, and a Riccati--type inequality for the smallest eigenvalue of
the Hessian. These structures underpin the curvature--flow mechanism
for mass generation in the Yang--Mills mass gap program.
\end{abstract}

\section{vHJ Flow for Densities Solving the Heat Equation}

Let $\rho_t:\mathbb{R}^n\to(0,\infty)$ be a smooth family of probability
densities solving the heat equation
\begin{equation}
  \partial_t\rho_t = \Delta\rho_t,
  \qquad \int_{\mathbb{R}^n}\rho_t(x)\,dx = 1.
\end{equation}

\begin{definition}
Assume there exist smooth functions $S_t:\mathbb{R}^n\to\mathbb{R}$ and
normalisation constants $Z_t>0$ such that
\begin{equation}
  \rho_t(x) = Z_t^{-1}e^{-S_t(x)}.
\end{equation}
We call $S_t$ the \emph{effective action} at time $t$.
\end{definition}

\begin{proposition}[Viscous Hamilton--Jacobi equation]
Up to an additive function of $t$ only, the effective action satisfies
\begin{equation}
  \partial_t S_t = \Delta S_t - |\nabla S_t|^2.
\end{equation}
\end{proposition}

\begin{proof}
Differentiate $\rho_t$:
\begin{equation}
  \partial_t\rho_t
  = -Z_t^{-1}(\partial_t S_t)e^{-S_t} - Z_t^{-2}(\partial_t Z_t)e^{-S_t}
  = -(\partial_t S_t + \partial_t\log Z_t)\rho_t.
\end{equation}
On the other hand,
\begin{equation}
  \Delta\rho_t
  = \Delta(e^{-S_t})
  = e^{-S_t}\big(-\Delta S_t + |\nabla S_t|^2\big)
  = \rho_t\big(-\Delta S_t + |\nabla S_t|^2\big).
\end{equation}
Equating $\partial_t\rho_t$ and $\Delta\rho_t$ and dividing by $\rho_t$
gives
\begin{equation}
  -\partial_t S_t - \partial_t\log Z_t
  = -\Delta S_t + |\nabla S_t|^2.
\end{equation}
The term $\partial_t\log Z_t$ depends only on $t$; absorbing it into
$S_t$ by redefining $\tilde S_t = S_t + \log Z_t$ leaves $\rho_t$ unchanged
and yields
\begin{equation}
  \partial_t \tilde S_t = \Delta \tilde S_t - |\nabla\tilde S_t|^2.
\end{equation}
Dropping the tilde gives the claim.
\end{proof}

\section{Hessian Flow and Reaction--Diffusion Structure}

Let $b_t = \nabla S_t$ and $h_t = \nabla^2 S_t$ denote the gradient and
Hessian of $S_t$. Differentiating the vHJ equation gives evolution
equations for $b_t$ and $h_t$.

\begin{proposition}[Gradient evolution]
The gradient evolves according to
\begin{equation}
  \partial_t b_t = \Delta b_t - 2(\nabla S_t\cdot\nabla)b_t.
\end{equation}
\end{proposition}

\begin{proof}
Differentiate $\partial_t S_t = \Delta S_t - |b_t|^2$:
\begin{align}
  \partial_t b_{t,i}
  &= \partial_i\partial_t S_t
   = \partial_i\Delta S_t - \partial_i\big(|b_t|^2\big) \nonumber\\
  &= \Delta b_{t,i} - 2\sum_j b_{t,j}\partial_i b_{t,j}
   = \Delta b_{t,i} - 2(\nabla S_t\cdot\nabla)b_{t,i}.
\end{align}
\end{proof}

\begin{proposition}[Hessian evolution]
The Hessian evolves according to
\begin{equation}
  \partial_t h_t
  = \Delta h_t - 2(\nabla S_t\cdot\nabla)h_t - 2h_t^2,
\end{equation}
where $(h_t^2)_{ik} = \sum_j h_{t,ij}h_{t,jk}$.
\end{proposition}

\begin{proof}
Differentiate the gradient equation:
\begin{align}
  \partial_t h_{t,ik}
  &= \partial_k\partial_t b_{t,i}
   = \partial_k\Delta b_{t,i}
     - 2\partial_k\Big(\sum_j b_{t,j}\partial_i b_{t,j}\Big) \nonumber\\
  &= \Delta h_{t,ik}
     - 2\sum_j\big[(\partial_k b_{t,j})(\partial_i b_{t,j})
                    + b_{t,j}\partial_i\partial_k b_{t,j}\big].
\end{align}
Now $\partial_k b_{t,j} = h_{t,jk}$ and
$\partial_i\partial_k b_{t,j} = \partial_i h_{t,jk}$, so
\begin{equation}
  \partial_t h_{t,ik}
  = \Delta h_{t,ik}
     - 2\sum_j\big(h_{t,jk}h_{t,ij} + b_{t,j}\partial_i h_{t,jk}\big).
\end{equation}
In matrix notation the first term inside the sum is $(h_t^2)_{ik}$ and
the second term is $(\nabla S_t\cdot\nabla h_t)_{ik}$. This yields
\begin{equation}
  \partial_t h_t
  = \Delta h_t - 2(\nabla S_t\cdot\nabla)h_t - 2h_t^2.
\end{equation}
\end{proof}

Thus each eigenvalue of the Hessian evolves under a combination of
diffusion, advection along the gradient flow, and a local quadratic
reaction term $-2\lambda^2$ that tends to drive curvature towards zero.

\section{Riccati Inequality for the Smallest Eigenvalue}

Let $\lambda(t,x)$ denote the smallest eigenvalue of $h_t(x)$ at the
point $x\in\mathbb{R}^n$. We give a local differential inequality for
$\lambda$ that has Riccati form. For simplicity we ignore geometric
subtleties and state the result in a formal, yet standard, way.

\begin{theorem}[Riccati inequality]
Assume $S_t$ and its derivatives are smooth and globally bounded so that
all the above operations are justified. Then at any point where $\lambda$
is attained by a simple eigenvalue with unit eigenvector $v$, one has
\begin{equation}
  \partial_t\lambda \ge \Delta\lambda - 2(\nabla S_t\cdot\nabla)\lambda
    - 2\lambda^2.
\end{equation}
Formally, $\lambda$ therefore satisfies the scalar parabolic inequality
\begin{equation}
  \partial_t\lambda \;\gtrsim\; \mathcal{L}_t\lambda - 2\lambda^2,
\end{equation}
where $\mathcal{L}_t = \Delta - 2\nabla S_t\cdot\nabla$ is the generator
of the Langevin diffusion associated with $S_t$.
\end{theorem}

\begin{proof}[Sketch]
Let $v$ be a normalised eigenvector of $h_t(x)$ with eigenvalue
$\lambda$. Differentiating the Rayleigh quotient
$\lambda = \langle v,h_tv\rangle$ in time and using the evolution
equation for $h_t$, together with suitable choice of local coordinates,
yields
\begin{equation}
  \partial_t\lambda
  = \langle v,\partial_t h_t v\rangle
  = \langle v,\Delta h_t v\rangle - 2\langle v,h_t^2 v\rangle
    - 2\langle v,(\nabla S_t\cdot\nabla h_t)v\rangle.
\end{equation}
The diffusion term $\langle v,\Delta h_t v\rangle$ can be written as
$\Delta\lambda$ plus non--negative contributions coming from variations
of $v$; similarly, the advection term produces $-2(\nabla S_t\cdot\nabla)
\lambda$ plus higher--order corrections. The quadratic term is
$\langle v,h_t^2 v\rangle = \lambda^2$ plus a non--negative remainder.
Collecting terms and discarding the non--negative remainders yields the
inequality.
\end{proof}

If we now \emph{drop} the diffusion and advection terms (which can only
help $\lambda$ grow, heuristically) we obtain the purely temporal
Riccati ODE
\begin{equation}
  \frac{d\ell}{dt} = -2\ell^2 + \sigma(t),
\end{equation}
where $\sigma(t)$ represents any additional positive source term (for
instance from intrinsic curvature or anomaly contributions in a more
general setting).

\begin{lemma}[Riccati lower bound]
Suppose $\ell$ solves
\begin{equation}
  \dot\ell = -2\ell^2 + \sigma_*,\qquad \ell(0)\ge \ell_0>0,
\end{equation}
with a constant $\sigma_*>0$. Then $\ell(t)$ remains bounded away from
zero and converges monotonically to the fixed point
$\ell_* = \sqrt{\sigma_*/2}$ as $t\to\infty$.
\end{lemma}

\begin{proof}
The ODE is explicitly solvable, but monotonicity and convergence can be
seen directly. The right--hand side is positive whenever
$\ell^2 < \sigma_*/2$ and negative when $\ell^2>\sigma_*/2$. Thus
$\ell(t)$ is driven towards the interval $[-\sqrt{\sigma_*/2},
\sqrt{\sigma_*/2}]$. Since we start with $\ell(0)>0$, the solution is
monotonically increasing up to $\sqrt{\sigma_*/2}$ if $\ell(0)$ is
smaller, or decreasing down to it if larger. Standard comparison
arguments show convergence.
\end{proof}

\section{Interpretation for Curvature--Driven Mass Generation}

In applications to field theory, $S_t$ is an effective action on an
(infinite--dimensional) configuration space, $\lambda(t)$ is the
smallest \emph{horizontal} eigenvalue of the Hessian of $S_t$, and
$\sigma(t)$ is supplied by a positive anomaly source (for example, the
trace anomaly). The parabolic comparison principle suggests that
$\lambda(t)$ is bounded below by the solution of a Riccati equation with
a strictly positive source:
\begin{equation}
  \dot\lambda_{\min}(t) \gtrsim -\alpha\lambda_{\min}(t)^2 + \sigma_*,
\end{equation}
with $\alpha>0$ capturing geometric constants. The lemma then shows that
$\lambda_{\min}(t)$ cannot decay to zero, but rather stabilises to a
strictly positive value of order $\sqrt{\sigma_*/\alpha}$.

Via the Bakry--\'Emery theory, a uniform positive lower bound on the
Hessian implies positive curvature in the sense of $\Gamma_2$ and hence
Poincar\'e and log--Sobolev inequalities with positive constants. These
translate into a spectral gap for the associated Langevin generator and
hence, in a suitable OS reconstruction, into a physical mass gap.

\begin{remark}
On a lattice configuration space $SU(N)^{|\mathcal{B}|}$, the Haar
measure provides a time--independent, strictly positive contribution to
$\sigma_*$, while the Wilson action and renormalisation effects contribute
through $\alpha$ and additional corrections to $\sigma(t)$. The finite--dimensional
derivations here can be lifted to that setting with only geometric
modifications (Lichnerowicz Laplacians and curvature tensors).
\end{remark}

\end{document}
